% ============================================================
%  Заключение
% ============================================================
\chapter*{Заключение}
\addcontentsline{toc}{chapter}{Заключение}

В работе систематизированы основные классы экзотических деривативов,
выплата по которым зависит от траектории~--- азиатские, барьерные,
радужные опционы и корреляционные свопы~--- и обозначены границы
применимости известных аналитических методов их оценки. Подробно
рассмотрены корреляционные свопы: приведены модели динамики
корреляции (constant, mean-reverting, Якоби, Фишер--OU) и выведены
выражения для справедливого страйка. Описано свойство аддитивности
forward-стоимости свопа на интервалы $[0,t]$ и $[t,T]$, лежащее в
основе динамического пересмотра позиции.

Для rainbow-опционов предложена декомпозиция цены на аналитически
считаемые компоненты (форвард, ванильные опционы, спред-опцион) и
оста\-точную нелинейную часть. Сравнены статическая и динамическая
стратегии хеджирования; при моделировании транзакционных издержек
показана нелинейная зависимость качества хеджа от частоты
ребалансировки. Сделан акцент на ключевой идее <<цена~=~стоимость
хеджа>>: следует устранять всё, что можно устранить дельта-хеджем и
явным <<выбрасыванием>> волатильностной и корреляционной экспозиции, а
обучаться~--- лишь на остатке, который проще и стабильнее
аппроксимировать нейронной сетью, чем подгонять модель сразу
к итоговой цене инструмента.

Эмпирическая часть демонстрирует, что подход с декомпозицией и
обучением остатка $R$ даёт более низкие значения RMSE и MAE по
сравнению с прямой ML-моделью на цену rainbow-опциона.

\todoblock{После проведения экспериментов уточнить числовые выводы и
добавить замечания по перспективам: расширение на $n\ge 3$ активов,
интеграция с моделями стохастической волатильности и корреляции,
оценка греков через автодифференцирование.}

\newpage
